%! Scatterplots and Correlation — Unit 2, Lesson 2.2 — Lesson Plan
% GRADUAL-RELEASE plan at the COMPACT budget (2026-09-29): warm-up / guided notes and practice /
% debrief / close and start the homework, 5 / 34 / 8 / 8. Teacher-facing.
% THERE IS NO GROUP ACTIVITY and NO INDEPENDENT PRACTICE SET: the notes are one Main Ideas /
% Notes table (four rows, 13 problems), the last row Guided Practice; the HOMEWORK is the
% individual practice, started in class. Every box is `skillbox` with a \boxguard ahead of it.
\documentclass[10pt]{article}
\usepackage{apstats-article}
\usepackage{apstats-boxes}

\newcommand{\UnitNumberName}{Unit 2: Exploring Two-Variable Data \quad}
\newcommand{\LessonNumberName}{Lesson 2.2: Scatterplots and Correlation}

\begin{document}
\begin{center}
    {\Large\bfseries \CourseName \\
    {\small\bfseries \UnitNumberName \LessonNumberName}}
\end{center}

\begin{tcolorbox}[colback=sky, colframe=navy, boxrule=0.9pt, arc=2mm,
    left=2mm, right=2mm, top=1.5mm, bottom=1.5mm]
    \textbf{Primary Objective:} Students will read a scatterplot, name its explanatory and
    response variables, describe the association by direction, form, strength, and unusual
    features in context, and interpret a given correlation $r$ --- including why an $r$ near 0
    does not mean ``no relationship'' and why a strong $r$ does not show causation.
    \\[2pt]
    \textbf{CED alignment:} Topics 2.4, 2.5 \quad$\cdot$\quad Big Ideas: UNC, DAT
    \quad$\cdot$\quad Skills 2.A, 2.B, 2.C, 4.B \quad$\cdot$\quad LO UNC-1.S, DAT-1.A,
    DAT-1.B, DAT-1.C \quad$\cdot$\quad EK UNC-1.S.1--1.S.3, DAT-1.A.1--1.A.6, DAT-1.B.1--1.B.3,
    DAT-1.C.1--1.C.2
    \\[2pt]
    \textbf{Lesson model:} \emph{Gradual release --- I do, we do, you do.} There is no group
    activity. The guided notes deliver the vocabulary and the core idea directly, row by row
    in one \emph{Main Ideas / Notes} table --- each idea a definition to complete and a grid of
    short problems, the first worked by you and the rest by the class; the last row, Guided
    Practice, is worked together with students holding the pen;
    \textbf{the homework is the individual practice}, and students start it in class while you
    circulate. The whole-class debrief is \emph{spoken}.
    \\[2pt]
    \textbf{Scope:} students \emph{read} $r$ from technology output and interpret it; they do
    not compute it by hand (the formula in EK DAT-1.B.2 is named, not worked). Every scatterplot
    is pre-drawn --- no one plots points today.
\end{tcolorbox}

\renewcommand{\arraystretch}{1.4}
\boxguard
\begin{skillbox}[Learning Targets \& Key Understandings]{goldbox}
\begin{tabularx}{\linewidth}{@{}X|X@{}}
\textbf{Learning targets (student language)} & \textbf{Key understandings (the ``why'')} \\[2pt]
\hline
\vspace{-6pt}
\begin{itemize}[leftmargin=1.1em, nosep, after=\vspace{2pt}]
  \item I can read a scatterplot and say which variable explains or predicts the other.
  \item I can describe an association by its direction, form, strength, and unusual points,
        in the words of the problem.
  \item I can say what a correlation tells me about a linear pattern --- and what it cannot.
  \item I can explain why two variables that move together may not cause each other.
\end{itemize}
&
\vspace{-6pt}
\begin{itemize}[leftmargin=1.1em, nosep, after=\vspace{2pt}]
  \item A scatterplot shows two quantitative variables on the same individuals; the
        explanatory variable goes on $x$ (EK UNC-1.S.1--1.S.3).
  \item A description names direction, form, strength, and unusual features; unusual means
        far from the \emph{pattern} (EK DAT-1.A.1--1.A.6).
  \item $r$ gives the direction and strength of the \emph{linear} association; it lies in
        $[-1, 1]$, has no units, and ignores which variable is $x$ (EK DAT-1.B.1, DAT-1.C.1).
  \item \textbf{Target misconception:} that $r$ near 0 means ``no relationship.'' One car
        driven at 12 speeds traces a clean arch and has $r = 0.03$: $r$ near 0 means no
        \emph{linear} association (notes problem~9). Its twin: that a strong $r$ proves
        causation (EK DAT-1.C.2; problems 10, 13).
\end{itemize}
\end{tabularx}
\end{skillbox}

\boxguard
\begin{skillbox}[{Vocabulary, Concepts, \& Theorems --- taught directly in the guided notes}]{greenbox}
\begin{minipage}[t]{0.48\linewidth}
    \begin{itemize}
        \item \textbf{Scatterplot} --- two quantitative variables, one dot per individual
              (EK UNC-1.S.2)
        \item \textbf{Explanatory \& response variables} --- $x$ explains or predicts $y$
              (EK UNC-1.S.3)
        \item \textbf{Association} --- direction (positive / negative), form (linear /
              curved), strength (strong / moderate / weak), unusual features --- DOFS
              (EK DAT-1.A)
    \end{itemize}
\end{minipage}
\hfill
\begin{minipage}[t]{0.48\linewidth}
    \begin{itemize}
        \item \textbf{Correlation $r$} --- direction and strength of a \emph{linear}
              association; $-1 \le r \le 1$; unit-free (EK DAT-1.B.1, DAT-1.C.1)
        \item \textbf{$r$ near $\pm1$} need not mean a line is the right model
              (EK DAT-1.B.3) --- returns in Lesson~2.5
        \item \textbf{Correlation is not causation} (EK DAT-1.C.2)
        \item \textbf{Carried forward} --- association between categorical variables (2.1);
              shape, center, spread, outliers (Unit~1)
    \end{itemize}
\end{minipage}
\end{skillbox}

% A tabularx never splits, so guard generously ahead of it.
\boxguard[30]
\begin{skillbox}[Lesson at a Glance --- \MeetingLength]{sky}
\renewcommand{\arraystretch}{1.25}
\begin{tabularx}{\linewidth}{@{}l c X X@{}}
\textbf{Phase} & \textbf{Min} & \textbf{Students} & \textbf{Teacher} \\
\hline
Warm-Up & 5 & Pass rates, the four words for a distribution, five students' study hours & Debrief item~2 aloud; leave ``shape, center, spread, outliers'' on the board \\
Guided Notes \& Practice & 34 & Fill the vocabulary box and the notes table: problems 1--10 row by row, then \emph{The City Pool} (11--13) together, pen in their hand & \textbf{I do} each row's definition lines and first problem (1, 4, 7), \textbf{we do} the rest of its grid (24) $\to$ \textbf{we do} Guided Practice (10) \\
Debrief & 8 & Pens down, papers up --- answer aloud, nothing to fill in & Open on the almost-right answer to problem~9; then causation \\
Close \& Assign & 8 & \textbf{Start the homework in class}, alone and silent & Announce the paper homework, circulate for your formative read on B2, preview Lesson~2.3 \\
\end{tabularx}
\end{skillbox}

\boxguard[20]
\begin{skillbox}[Warm-Up --- Activate Prior Knowledge (5 min)]{sky}
\begin{minipage}[t]{0.48\linewidth}
    \textbf{The three items and what each seeds}
    \begin{itemize}
        \item \textbf{1.} Pass rates of 75\% vs.\ 60\%. \textbf{Spirals} to Lesson 2.1 and plants
              the word \emph{association}: knowing one variable helps predict the other.
        \item \textbf{2.} Shape, center, spread, outliers. \textbf{Seeds} DOFS --- a scatterplot
              gets its own four words in row~2.
        \item \textbf{3.} Five study-hours / score pairs. \textbf{Seeds} direction and
              explanatory vs.\ response before either has a name.
    \end{itemize}
\end{minipage}
\hfill
\begin{minipage}[t]{0.48\linewidth}
    \textbf{Running it}
    \begin{itemize}
        \item \textbf{Debrief item~2 aloud} and leave the four words on the board; in row~2
              write DOFS beside them so students see the parallel.
        \item Item~3's second question has a defensible answer either way; take ``hours,
              because it comes first'' and move on --- row~1 names it.
        \item If item~1 stalls, give the two percents and ask only \emph{``do they differ?''}
    \end{itemize}
\end{minipage}
\end{skillbox}

\boxguard[30]
\begin{skillbox}[Guided Notes \& Practice --- I do, then we do (34 min)]{sky}
\textbf{This is the whole lesson.} There is no group activity, and there is no independent
practice set on this page: \textbf{the homework is the individual practice}, started in class.
\begin{multicols}{2}
The notes are one \emph{Main Ideas / Notes} table --- four rows, 13 short problems.
\textbf{In every row you fill the definition lines and work the first problem of the grid; the
class works the rest, pen in their hand.} The vocabulary box (three terms) is filled \emph{as
you go}, never in advance.

\textbf{Variables (5 min).} Name \emph{scatterplot} and \emph{explanatory/response} on the 12
cars. Work problem~1 yourself; 2 and 3 are theirs and quick --- one dot is one \emph{car model},
not one gallon.

\textbf{Pattern (8 min).} DOFS beside the warm-up's four words. Work problem~4 (direction, in
context). Problem~5 is theirs. \textbf{Problem~6 is the trap}: the heaviest car is far to the
right but sits on the line; the hybrid at 48\,mpg is the unusual point. Take a hand count
before anyone speaks.

\textbf{Reading $r$ (11 min) --- must-land.} Fill the definition, then work problem~7
($r = -0.86$: sign and size). Problem~8 is fast --- units and axis order do not move $r$.
\textbf{Problem~9 is the crux}: put the speed plot up, say ``$r = 0.03$,'' and let a student
say ``no relationship'' before anyone looks. Then look. Problem~10 is causation; accept
``national wealth'' or anything that drives both.

\columnbreak
\textbf{We do (about 10 min): The City Pool.} Problems 11--13, a fresh context, worked
entirely together. Students hold the pen; you hold the questions: \emph{``Which one would you
change to see the other move?''} (11) \quad \emph{``Give me all four words, then $r$, in pool
language''} (12) \quad \emph{``What else goes up on a hot day?''} (13). Problem~12 is the one to
protect --- it is the full description the AP exam scores.

\textbf{The pen must actually change hands.} Guided Practice is the only place today where you
can watch a student decide, so do not narrate it. Circulate and demand the reason, not the
answer:
\begin{itemize}[leftmargin=1.1em, nosep]
  \item \textbf{Problem~12:} ``You said strong. Strong \emph{what}? Say the form.''
  \item \textbf{Problem~13:} ``If the snack bar closed, would fewer people swim?''
\end{itemize}
While circulating, pick the papers for the debrief --- one problem~9 that says ``weak
relationship'' (almost right), and one that names the curve. \textbf{Your formative read comes
twice}: here, and again in the last eight minutes as students start the homework.
\end{multicols}
\end{skillbox}

\boxguard
\begin{skillbox}[Debrief --- whole class (8 min)]{sky}
Pens down, papers up. \textbf{This debrief is spoken} --- nothing on the page to fill in.
\begin{multicols}{2}
\begin{enumerate}[leftmargin=1.3em, nosep, itemsep=3pt]
  \item \textbf{Open on the \emph{almost}-right answer to problem~9} --- ``a weak relationship.''
        Ask: \emph{``Weak? If I tell you the speed, how well can you guess the mpg?''} Very
        well. So the relationship is strong; only the \emph{linear} part is missing. Land
        the sentence: $r$ near 0 means no linear association.
  \item Put $r = -0.86$ (cars) beside $r = 0.03$ (speed) and ask which plot is easier to
        predict from. The answer is ``it depends on the shape'' --- look first.
  \item Take problems 10 and 13 together: in each, name the third variable that drives both.
\end{enumerate}
\columnbreak
\textbf{Two things to demand aloud.} A student states what $r$ \emph{does} tell you (direction
and strength of a linear pattern), and another states one thing only the picture can tell you
(a curve, a cluster, an unusual point).

\textbf{If the debrief runs short}, cut item~2 --- item~1 is the only one that cannot be cut.
\textbf{Do not let the debrief eat the last eight minutes} --- the homework start is your
second formative read.
\end{multicols}
\end{skillbox}

\boxguard
\begin{skillbox}[AP Practice --- assigned, not scored]{goldbox}
One page on 20 teenagers' screen time and sleep ($r = -0.72$). \textbf{MC1} interprets $r$
(distractors: ``72\% caused by,'' ``slope of 0.72,'' ``negative means weak''). \textbf{MC2}
asks which change moves $r$: units, axis swap, and a constant shift do not; a new point off the
pattern does. \textbf{FRQ 3}: (a) variables and a full description in context; (b) a causal
headline to reject; (c) \textbf{the spiral} to Lesson 2.1 --- recorded as categories, the
association lives in a two-way table.

\textbf{Not scored --- the cover's score column reads \textbf{NA} for this page.} Go over it at
the start of Lesson~2.3 and hold the AP standard out loud: \emph{in context or it does not
count}. Part (a) earns the point only with all of direction, form, and strength \emph{and}
the variable names; (b) must name a plausible other cause, not just recite ``correlation is not
causation.''
\end{skillbox}

\boxguard
\begin{skillbox}[Homework --- scored, due the first class after two study halls]{goldbox}
Five items in a third context. \textbf{Part A} (a basketball team, practice hours vs.\
free-throw \%, $r = 0.60$): \textbf{A1} explanatory and response; \textbf{A2} a full DOFS
description, including the 11-hour, 50\% player who sits far below the pattern. \textbf{Part B}
(what $r$ can and cannot tell you): \textbf{B1} an AP-style multiple choice --- strongest $r$
is $-0.91$ --- with a one-sentence justification; \textbf{B2} the target misconception, a
U-shaped age / reaction-time plot with $r = 0.05$; \textbf{B3} libraries and crime, with
population as the lurking variable.

\textbf{Start it in the last eight minutes of the period, in class and alone.}
\textbf{Scoring:} 10 points --- 2 per item. \textbf{B2 is the one to read closely}: it is
problem~9 without the picture, and a student who writes ``Jordan is right, $r$ is almost 0''
has not yet separated \emph{linear} from \emph{related}. Lesson~2.3 fits a line; a student who
cannot see that a curve is not a line will fit one anyway.

\textbf{Paper homework is the default for this course} --- assign this page unless you have
decided otherwise. \textbf{DeltaMath override (occasional):} on the days you assign DeltaMath
instead, it replaces this page, you strike through the score cell on the cover, and this page
stays in the packet as extra practice. Never assign both.
\end{skillbox}

\boxguard
\begin{skillbox}[Watch For (while circulating)]{redbox}
\begin{itemize}[leftmargin=1.2em, nosep]
  \item \textbf{``$r$ is near 0, so there is no relationship''} (notes~9, HW~B2 --- the target
        misconception). Probe: \emph{``If I tell you the speed, can you guess the mpg?''}
  \item \textbf{A strong $r$ read as cause} (notes~10, 13; AP~3(b); HW~B3). Probe:
        \emph{``What else goes up when both of these do?''}
  \item \textbf{``Negative means weak''} (AP~MC1, HW~B1). Probe: \emph{``Which is closer to
        the edge of the number line, $-0.91$ or $0.78$?''}
  \item \textbf{Extreme $x$ called an outlier} (notes~6, HW~A2). Probe: \emph{``Is it far from
        the others, or far from the pattern?''}
  \item \textbf{Descriptions with no context} --- ``strong negative linear'' with no variables
        (notes~12, AP~3(a), HW~A2). Probe: \emph{``Strong negative linear \emph{what}?''}
\end{itemize}
Cold-call prompts: \emph{``Give me a pair of variables with a strong relationship and $r$ near
0.''} \quad \emph{``What happens to $r$ if we measure in kilograms?''}
\end{skillbox}

\boxguard
\begin{skillbox}[Close \& Assign (8 min)]{goldbox}
\textbf{Students start the homework here, in class, alone and silent --- this is the individual
practice, and the first minutes of it are your best formative read.} Hand the launch first ---
five items on paper, scored, due the first class after two study halls --- and point out that
the AP Practice page before it is theirs to work and bring to review. Circulate and read
\textbf{B2} over shoulders; sort as you walk --- ``no, there is a curve, $r$ only measures
linear'' (ready); ``no'' with no reason (ready, needs the language); ``yes, $r$ is almost 0''
(the misconception survived) --- and open Lesson~2.3 by putting the speed plot back up if that
third pile ran past a third of the room. Then close on what changed: \emph{``You came in
thinking a number could tell you whether two things are related. You leave knowing $r$ only
measures a straight line --- look at the picture first.''} \textbf{Preview:} Lesson~2.3,
\emph{Regression Lines and Residuals} --- when the pattern is linear, we draw the line and
use it to predict.
\end{skillbox}

% ── Teacher notes — one per component, in packet order ──
\begin{teachernote}[Warm-Up]
Five minutes, and item~2 is the one to debrief: the four words for a distribution go on the
board and stay there, because row~2 of the notes hangs DOFS beside them. Item~1 is a 60-second
spiral to Lesson 2.1; if it drags, give the percents. Item~3 has no wrong answer to its second
question --- let a student say ``hours, because it comes first'' and promise that row~1 names it.
\end{teachernote}

\begin{teachernote}[Guided Notes \& Practice]
\textbf{Pace it 24 / 10, then 8 for the debrief, then 8 for the homework start.} In every row
you fill the lines and work the first problem (1, 4, 7); the class works the rest. Rows~1 and~2
are machinery and must move --- the cars carry both, so no time is spent loading a second
context. \textbf{Problem~6 is the trap}: students call the heaviest car the outlier because it
is far from the others; the hybrid is far from the \emph{pattern}.

\textbf{Spend the time in Reading $r$.} Say ``$r = 0.03$'' \emph{before} students look at the
speed plot and take the ``no relationship'' answer publicly --- problem~9 lands only if somebody
commits to the wrong answer first. Then point at the arch: best mileage near 45--50\,mph,
worse on either side, and the pattern is as tight as the cars'. Do not compute $r$; the formula
is technology's job (EK DAT-1.B.2).

\textbf{The City Pool is where the pen changes hands.} If you only get through 11 and 12, that
is the right trade --- 13 is causation again and the debrief covers it.

\textbf{There is no independent practice set, and that is deliberate --- the homework is the
individual practice.} Read HW~B2 during the supervised start and sort into the three piles in
the Close box. If the third pile is more than a third of the room, reopen the speed plot at the
start of Lesson~2.3. \textbf{Debrief (8 min), spoken} --- open on the ``weak relationship''
answer to problem~9; it is the one move that cannot be cut.
\end{teachernote}

\begin{teachernote}[AP Practice]
\textbf{This page is not required and not scored} --- the graded practice is the homework behind
it. Go over it at the start of Lesson~2.3. \textbf{MC2 is the item worth reading closely}: (A),
(B), and (C) are exactly the changes $r$ ignores, so a student who picks any of them has
memorized ``$r$ has no units'' without seeing why. FRQ~(b) earns credit only with a named
alternative cause (homework load, a late bedtime); FRQ~(c) is the spiral to Lesson 2.1.
\end{teachernote}

\begin{teachernote}[Homework]
\textbf{Scored, 10 points (2 per item), due the first class after two study halls.} The eight
minutes in class are a \emph{start}: put students on A1, A2, and B2 while you can still catch a
wrong start. \textbf{B2 predicts Lesson~2.3} --- a student who reads $r = 0.05$ as ``unrelated''
will fit a line to a curve next time. A2 earns full credit only in context and only if it
flags the 11-hour player; B1 needs the justification, not just the letter.
\textbf{When DeltaMath carries this topic, assign it instead}, strike the score cell on the
cover, and tell students this page is now optional practice. Do not assign both.
\end{teachernote}

\end{document}
